%% ma: L11-B14-0009
%% lop: 11
%% chuong: 4
%% bai: 14
%% dang: 11.14.2
%% loai: TN4
%% muc: TH
%% dap_an: "A"
%% nguon: "tai-lieu-hoc-tap-mon-toan-lop-11-theo-chuong-trinh-giao-duc-pho-thong-2018.pdf (Lớp Toán Thầy Duy, Phạm Lê Duy, THPT Chu Văn An), Chương 4 Bài 5, Phần 1, Câu 4 (tr.526)"
%% kien_thuc: "Lớp 11 Bài 14 (phép chiếu song song)"
%% trang_thai: chinh-thuc
%% ghi_chu: "Đã duyệt 03/11/2026."
\begin{ex}
Trong các mệnh đề sau, mệnh đề nào đúng?
\choice{\True Hình chiếu song song của hai đường thẳng chéo nhau có thể song song với nhau.}{Hình chiếu song song của hai đường thẳng cắt nhau thì song song.}{Hình chiếu song song của một hình vuông là một hình vuông.}{Hình chiếu song song của một lục giác đều là một lục giác đều.}
\loigiai{\sloppy A đúng: ví dụ hai đường chéo nhau $a,b$ cùng nằm trong hai mặt phẳng song song với phương chiếu thì hình chiếu của chúng song song. B sai: hai đường cắt nhau có hình chiếu cắt nhau (hoặc trùng nhau). C, D sai: hình chiếu của hình vuông là hình bình hành, của lục giác đều là lục giác có các cặp cạnh đối song song, không nhất thiết là hình vuông hay lục giác đều.}
\end{ex}
